En una clase de geometría, un estudiante analiza un triángulo y observa un punto marcado dentro de la figura, el cual corresponde a la intersección de las mediatrices de sus lados.

\begin{center}
\begin{tikzpicture}[scale=0.8]

% Triángulo
\coordinate (A) at (0,0);
\coordinate (B) at (6,0);
\coordinate (C) at (2,4);

\draw (A)--(B)--(C)--cycle;

% Etiquetas
\node[below left] at (A) {A};
\node[below right] at (B) {B};
\node[above] at (C) {C};

% Circuncentro aproximado
\coordinate (O) at (3,1.5);
\fill (O) circle (2pt);
\node[right] at (O) {$O$};

% Mediatriz vertical
\draw[dashed] (3,-1) -- (3,3);

% Marca de ángulo recto
\draw (3,0) -- (3,0.3) -- (3.3,0.3) -- (3.3,0);

\end{tikzpicture}
\end{center}

Se afirma que el punto $O$ es el circuncentro del triángulo.

Para comprobar esta afirmación, se propone el siguiente procedimiento con compás.

¿Cuál de las siguientes acciones permite verificar correctamente que $O$ es el circuncentro?

\begin{multicols}{2}
\begin{enumerate}[label=\alph*), leftmargin=*]
\item Trazar una circunferencia con centro en $O$ que pase por un vértice y verificar que también pasa por los otros dos vértices
\item Medir los lados del triángulo y comprobar que todos son iguales
\item Trazar segmentos desde $O$ a los lados del triángulo y verificar que son perpendiculares
\item Dibujar una circunferencia con centro en un vértice del triángulo
\end{enumerate}
\end{multicols}